
\documentclass[12pt, a4paper]{article}

\usepackage[margin=1.2in]{geometry}
\usepackage{booktabs}
\usepackage{hyperref}
\usepackage{enumitem}
\usepackage{parskip}
\usepackage{setspace}
\usepackage{xcolor}
\usepackage{listings}
\usepackage{titlesec}
\usepackage{fancyhdr}
\usepackage{amsmath}
\usepackage{tikz}
\usetikzlibrary{shapes.geometric, arrows.meta, positioning, fit}

\onehalfspacing

\hypersetup{colorlinks=true, linkcolor=black, citecolor=black, urlcolor=blue}

\setlength{\headheight}{13.6pt}
\pagestyle{fancy}
\fancyhf{}
\rhead{\small Pipeline Documentation}
\lhead{\small Bank Fragility Project}
\cfoot{\thepage}
\renewcommand{\headrulewidth}{0.4pt}

\titleformat{\section}{\large\bfseries}{\thesection.}{0.5em}{}
\titleformat{\subsection}{\normalsize\bfseries}{\thesubsection}{0.5em}{}

\definecolor{codegray}{gray}{0.95}
\definecolor{darkblue}{RGB}{31, 56, 100}
\lstset{
    backgroundcolor=\color{codegray},
    basicstyle=\ttfamily\footnotesize,
    breaklines=true,
    frame=single,
    framerule=0pt,
}

\begin{document}

\begin{titlepage}
    \centering
    \vspace*{2cm}
    {\LARGE\bfseries Bank Fragility Analysis\\[0.4em]
    Pipeline Documentation \par}
    \vspace{1cm}
    {\large\itshape From Raw Data to Replicable Results \par}
    \vspace{2cm}
    {\normalsize
    \textbf{Report Date:} 12/31/2025 \\[0.4em]
    \textbf{Shock Window:} 2020-01-01 to 2025-12-31 \\[0.4em]
    \textbf{RMBS Multiplier:} 1.25
    \par}
    \vfill
    {\small All data sources are publicly available. No proprietary
    subscription required.\par}
\end{titlepage}

\tableofcontents
\newpage

%% ─────────────────────────────────────────────────────────────────────────────
\section{Overview}

This document describes the computational pipeline for the Bank Fragility
replication project. The pipeline collects, processes, and analyses U.S.\ bank
balance sheet data to replicate the fragility measures of Jiang, Matvos,
Piskorski, and Seru (2023).

The pipeline is automated using \texttt{doit}, a Python task runner that tracks
file dependencies and re-executes only tasks whose inputs have changed. Running
\texttt{doit} from the project root executes the full pipeline end-to-end.

\subsection*{Quick Start}

\begin{lstlisting}
# Install dependencies
pip install doit python-dotenv pandas yfinance selenium requests
pip install openpyxl scipy nbconvert

# Configure
cp .env.example .env   # edit REPORT_DATE, MARKET_START_DATE, etc.

# Run full pipeline
doit

# Run a single task
doit pull_treasury
doit make_table_1
\end{lstlisting}

%% ─────────────────────────────────────────────────────────────────────────────
\section{Dependency Graph}

\begin{center}
\begin{tikzpicture}[
    node distance=0.5cm and 2.0cm,
    box/.style={rectangle, draw, rounded corners=3pt, minimum width=3cm,
                minimum height=0.75cm, align=center, font=\small},
    pull/.style={box, fill=blue!12},
    proc/.style={box, fill=orange!18},
    out/.style={box, fill=green!15},
    arr/.style={-{Stealth[length=5pt]}, thick, gray},
]

\node[pull] (ffiec)   {pull\_ffiec};
\node[pull, below=of ffiec]  (gsib)    {pull\_gsib};
\node[pull, below=of gsib]   (tsy)     {pull\_treasury};
\node[pull, below=of tsy]    (mbs)     {pull\_mbs};

\node[proc, right=2.5cm of ffiec, yshift=-0.6cm] (proc)  {process\_ffiec};
\node[proc, right=2.5cm of tsy,   yshift=-0.6cm] (shock) {compute\_yield\\shocks};

\node[out,  right=2.5cm of proc,  yshift=-0.6cm] (tbl)   {make\_table\_1};
\node[out,  below=0.8cm of tbl]                  (nb)    {convert\_notebook};
\node[out,  below=0.8cm of nb]                   (pdf)   {generate\_pdf};

\draw[arr] (ffiec) -- (proc);
\draw[arr] (gsib)  -- (proc);
\draw[arr] (tsy)   -- (shock);
\draw[arr] (proc)  -- (tbl);
\draw[arr] (shock) -- (tbl);

\end{tikzpicture}
\end{center}

\begin{center}
\small
\textcolor{blue!60}{\rule{1cm}{4pt}} Data collection \quad
\textcolor{orange!70}{\rule{1cm}{4pt}} Processing \quad
\textcolor{green!60}{\rule{1cm}{4pt}} Analysis \& Output
\end{center}

%% ─────────────────────────────────────────────────────────────────────────────
\section{Data Collection Scripts}

\subsection{\texttt{pull\_ffiec.py} --- FFIEC Call Reports}

\textbf{Purpose:} Downloads the FFIEC Call Report bulk data archive for the
configured reporting quarter.

\textbf{Method:} Selenium WebDriver automates the FFIEC bulk download portal.
The script selects ``Call Reports -- Single Period'', chooses the reporting date,
ensures Tab-Delimited format is selected, and triggers the download. If the
requested date is unavailable, the most recent available date is used.

\textbf{Configuration:}
\begin{lstlisting}
REPORT_DATE_SLASH = 12/31/2025   # date shown on FFIEC website
REPORT_DATE       = 12312025         # date in zip filename
\end{lstlisting}

\textbf{Output:}
\begin{lstlisting}
_data/FFIEC CDR Call Bulk All Schedules 12312025.zip
\end{lstlisting}

\textbf{Notes:} Requires Google Chrome and ChromeDriver. The download may take
several minutes as the ZIP file is typically 500MB--1GB. The 15-second wait
after clicking download can be increased for slower connections.

\subsection{\texttt{pull\_gsib\_banks.py} --- GSIB Bank List}

\textbf{Purpose:} Generates a lookup table of Globally Systemically Important
Banks identified by their RSSD IDs.

\textbf{Method:} The list of 37 RSSD IDs is stored as a static constant in the
script. This reflects the set of institutions that have appeared on the FSB
GSIB list. The list is saved as a Parquet file for fast joins downstream.

\textbf{Output:}
\begin{lstlisting}
_data/gsib_list.parquet   # columns: rssd_id_call, is_gsib
\end{lstlisting}

\subsection{\texttt{pull\_treasury\_yields.py} --- Treasury Yields}

\textbf{Purpose:} Downloads the full history of U.S.\ Treasury constant
maturity yields from the FRED API.

\textbf{Method:} Six yield series are pulled via authenticated HTTP requests
to the FRED API endpoint. Each series is requested from 1900-01-01 to obtain
the full available history. Failed requests are retried up to three times with
a 2-second pause. FRED returns missing values as the string ``.'', which are
coerced to NaN and dropped.

\begin{center}
\begin{tabular}{lll}
\toprule
Series & Column & Coverage \\
\midrule
DGS1  & \texttt{dgs1}  & 1962--present \\
DGS3  & \texttt{dgs3}  & 1962--present \\
DGS5  & \texttt{dgs5}  & 1962--present \\
DGS10 & \texttt{dgs10} & 1962--present \\
DGS20 & \texttt{dgs20} & 1962--present (gap 1987--1993) \\
DGS30 & \texttt{dgs30} & 1977--present \\
\bottomrule
\end{tabular}
\end{center}

\textbf{Configuration:}
\begin{lstlisting}
FRED_API_KEY = <your key>   # free at fred.stlouisfed.org
\end{lstlisting}

\textbf{Output:}
\begin{lstlisting}
_data/treasury_yields.parquet   # one row per trading day
\end{lstlisting}

\subsection{\texttt{pull\_mbs\_etfs.py} --- MBS ETF Prices}

\textbf{Purpose:} Downloads daily price history for two MBS ETFs as market
proxies for residential and commercial mortgage-backed securities.

\textbf{Method:} The \texttt{yfinance} library is used to download adjusted
close prices from Yahoo Finance. Daily returns are computed as percentage
changes. A validation check confirms that the configured shock window is fully
covered by the downloaded data.

\begin{center}
\begin{tabular}{lll}
\toprule
Ticker & Columns & Description \\
\midrule
SPMB & \texttt{rmbs\_px}, \texttt{rmbs\_ret} & Residential MBS ETF \\
CMBS & \texttt{cmbs\_px}, \texttt{cmbs\_ret} & Commercial MBS ETF  \\
\bottomrule
\end{tabular}
\end{center}

\textbf{Output:}
\begin{lstlisting}
_data/mbs_etfs.parquet   # columns: date, rmbs_px, cmbs_px, rmbs_ret, cmbs_ret
\end{lstlisting}

%% ─────────────────────────────────────────────────────────────────────────────
\section{Processing Scripts}

\subsection{\texttt{process\_ffiec.py} --- FFIEC Data Processing}

\textbf{Purpose:} Extracts, cleans, and structures the raw FFIEC Call Report
data into a bank-level panel suitable for analysis.

\textbf{Steps:}
\begin{enumerate}[noitemsep]
    \item Extract five schedule files from the ZIP archive (RC, RCA, RCB, RCCI, RCE)
    \item Parse each as tab-delimited, standardise column names, convert to numeric
    \item Separate columns into \texttt{rcfd} (global/consolidated) and
          \texttt{rcon} (domestic) prefixes
    \item Construct asset variables: cash, securities by type, loans by type,
          fed funds, repo
    \item Allocate assets into six maturity buckets using reported breakdowns
          (RMBS) or stylized weights (other categories)
    \item Construct liability variables: deposits (insured/uninsured), fed funds
          purchased, repo, equity
    \item Merge global and domestic tables, filling gaps with available data
    \item Tag each bank as Small / Large / GSIB
    \item Compute summary statistics with winsorization at the 5th/95th percentiles
    \item Export bank panel, Excel workbook, figure, and LaTeX tables
\end{enumerate}

\textbf{Outputs:}
\begin{lstlisting}
_data/bank_panel_12312025.parquet
_output/summary_stats_12312025.xlsx
_output/figure_A1_12312025.png
_output/summary_assets.tex
_output/summary_liabilities.tex
\end{lstlisting}

\subsection{\texttt{compute\_yield\_shocks.py} --- Market Shock Parameters}

\textbf{Purpose:} Derives empirical interest rate shocks from actual Treasury
yield movements over the configured shock window.

\textbf{Method:} For each of six maturity buckets, the yield at the bucket
midpoint is interpolated linearly from the nearest two FRED maturities.
The shock is the change in this interpolated yield between the start and end
of the shock window.

\begin{center}
\begin{tabular}{lrp{5.5cm}}
\toprule
Bucket & Midpoint & Interpolated from \\
\midrule
\texttt{lt1y}    & 0.5 yr  & DGS1 \\
\texttt{1\_3y}   & 2.0 yr  & 0.5$\times$DGS1 + 0.5$\times$DGS3 \\
\texttt{3\_5y}   & 4.0 yr  & 0.5$\times$DGS3 + 0.5$\times$DGS5 \\
\texttt{5\_10y}  & 7.5 yr  & 0.5$\times$DGS5 + 0.5$\times$DGS10 \\
\texttt{10\_15y} & 12.5 yr & 0.75$\times$DGS10 + 0.25$\times$DGS20 \\
\texttt{15plus}  & 22.0 yr & 0.8$\times$DGS20 + 0.2$\times$DGS30 \\
\bottomrule
\end{tabular}
\end{center}

\textbf{Configuration:}
\begin{lstlisting}
MARKET_START_DATE = 2020-01-01
MARKET_END_DATE   = 2025-12-31
RMBS_MULTIPLIER   = 1.25
\end{lstlisting}

\textbf{Output:}
\begin{lstlisting}
_data/market_shocks.parquet   # one row: d_tsy_* columns + rmbs_multiplier
\end{lstlisting}

\subsection{\texttt{make\_table\_1.py} --- Bank Fragility Table}

\textbf{Purpose:} Computes bank-level mark-to-market losses and the primary
fragility measure, and aggregates results into Table 1.

\textbf{Method:} For each bank and each maturity bucket, losses are computed as:
\[
    \text{Loss}_{i,k} = H_{i,k} \times \Delta y_k \times D_k \times m_k
\]
where $H_{i,k}$ is holdings, $\Delta y_k$ is the yield shock, $D_k$ is the
duration midpoint, and $m_k$ is the RMBS multiplier for mortgage assets.

The fragility measure is:
\[
    \text{Fragility}_i = \frac{\text{Uninsured Deposits}_i}{A_i - L_i}
\]

Results are reported for four groups: all banks, small, large non-GSIB, and GSIB.

\textbf{Outputs:}
\begin{lstlisting}
_output/table_1.csv
_output/table_1.tex
\end{lstlisting}

%% ─────────────────────────────────────────────────────────────────────────────
\section{Document Scripts}

\subsection{\texttt{convert\_to\_html.py} --- Notebook to HTML}

Converts \texttt{documents/eda.ipynb} to a self-contained HTML report using
\texttt{nbconvert}. The \texttt{--no-input} flag suppresses code cells so
only outputs and plots are shown, making the report suitable for sharing.

\textbf{Output:} \texttt{documents/eda.html}

\subsection{\texttt{generate\_pipeline\_doc.py} --- This Document}

Generates this pipeline documentation PDF from a Python string template,
embedding current configuration values from \texttt{.env}.

\subsection{\texttt{generate\_research\_paper.py} --- Research Paper}

Generates a brief research paper PDF focused on the analysis results,
including actual output tables loaded from \texttt{\_output/}.

%% ─────────────────────────────────────────────────────────────────────────────
\section{Configuration Reference}

All parameters are set in \texttt{.env} at the project root:

\begin{center}
\begin{tabular}{lll}
\toprule
Variable & Value & Used by \\
\midrule
\texttt{REPORT\_DATE\_SLASH} & 12/31/2025 & \texttt{pull\_ffiec.py} \\
\texttt{REPORT\_DATE}        & 12312025     & \texttt{process\_ffiec.py}, \texttt{make\_table\_1.py} \\
\texttt{MARKET\_START\_DATE} & 2020-01-01 & \texttt{compute\_yield\_shocks.py} \\
\texttt{MARKET\_END\_DATE}   & 2025-12-31   & \texttt{compute\_yield\_shocks.py} \\
\texttt{RMBS\_MULTIPLIER}    & 1.25   & \texttt{compute\_yield\_shocks.py}, \texttt{make\_table\_1.py} \\
\texttt{FRED\_API\_KEY}      & \textit{(hidden)} & \texttt{pull\_treasury\_yields.py} \\
\bottomrule
\end{tabular}
\end{center}

%% ─────────────────────────────────────────────────────────────────────────────
\section{File Structure}

\begin{lstlisting}
Bank_Fragility/
|-- .env                          # Configuration (never commit this)
|-- dodo.py                       # doit pipeline definition
|-- scripts/
|   |-- pull_ffiec.py             # Download FFIEC ZIP (Selenium)
|   |-- pull_gsib_banks.py        # GSIB bank list
|   |-- pull_treasury_yields.py   # Treasury yields from FRED
|   |-- pull_mbs_etfs.py          # MBS ETF prices from Yahoo Finance
|   |-- compute_yield_shocks.py   # Empirical shock parameters
|   |-- compute_market_shocks.py  # Hardcoded shock version (reference)
|   |-- process_ffiec.py          # Parse and structure FFIEC data
|   `-- make_table_1.py           # Bank fragility Table 1
|-- document_scripts/
|   |-- convert_to_html.py        # eda.ipynb -> eda.html
|   |-- generate_pipeline_doc.py  # This document
|   `-- generate_research_paper.py# Analysis research paper
|-- documents/
|   |-- eda.ipynb                 # EDA notebook
|   |-- eda.html                  # Rendered notebook (output)
|   |-- pipeline_documentation.pdf
|   `-- research_paper.pdf
|-- _data/                        # Downloaded and processed parquets
`-- _output/                      # Tables, figures, LaTeX exports
\end{lstlisting}

\end{document}
